\bookchapter{Kapitel 5: Das Smartphone besser verstehen lassen – Grundeinstellungen}

\chapterintro{Ein paar Minuten Einstellungen verbessern das Erlebnis enorm. Diese Einstellungen sind wie das „Stimmen“ des Smartphones, damit es besser zu deinen Gewohnheiten passt.}

\booksection{Eins. Ton und Vibration}

Lege dir einen Klingelton und Benachrichtigungston fest. Nutze den Stummschaltmodus (in Meetings) und den Nicht-stören-Modus (beim Schlafen) – die Welt wird sofort ruhiger.

\begin{itemize}
  \item \textbf{Klingelton}: Wähle einen, den du hörst und der nicht peinlich ist.
  \item \textbf{Vibrationsmodus}: Im Stummschaltmodus per Vibration aufmerksam machen.
  \item \textbf{Nicht-stören-Modus}: Ausnahmen festlegen, z. B. Anrufe aus Kontakten erlauben.
  \item \textbf{Lautstärke getrennt}: Medien-, Klingel- und Weckerlautstärke können separat eingestellt werden.
\end{itemize}

\booksection{Zwei. Anzeige und Helligkeit}

\begin{itemize}
  \item \textbf{Automatische Helligkeit}: Das Smartphone passt sich dem Umgebungslicht an – bequem.
  \item \textbf{Augenschutz-/Nachtmodus}: Filtert blaues Licht; abends ist das Smartphone angenehmer zu lesen.
  \item \textbf{Dunkelmodus}: Spart Strom (besonders effektiv bei OLED-Bildschirmen) und blendet nachts nicht. Viele Apps unterstützen ebenfalls den Dunkelmodus.
  \item \textbf{Bildschirm-Aktualisierungsrate}: Bei einigen Android-Geräten einstellbar – je höher (z. B. 120 Hz), desto flüssiger das Scrollen.
  \item \textbf{Auto-Sperre}: Auf 15 oder 30 Sekunden setzen – sicher und stromsparend.
  \item \textbf{Schriftgröße}: Schrift zu klein? Etwas größer stellen.
\end{itemize}

\booksection{Drei. Hintergrundbild, Theme und Startbildschirm}

Ein schönes Hintergrundbild (statisch oder animiert) verändert die Stimmung. Android-Smartphones haben oft einen Theme-Store, mit dem sich Icons und Stil komplett ändern lassen.

\textbf{Startbildschirm-Widgets}: Wetter, Kalender, Musiksteuerung direkt auf dem Startbildschirm – ein Schritt bis zur Aktion. iOS und Android unterstützen dies; lange auf eine freie Stelle des Startbildschirms drücken, um Widgets hinzuzufügen.

\booksection{Vier. Effizienzeinstellungen: Das Smartphone als dein „externes Gehirn“}

Dieser Bereich behandelt Funktionen, „die man nach dem Einrichten nicht mehr missen möchte“ – jede einzelne kann deine tägliche Effizienz spürbar steigern.

\booksubsection{Fokus / Nicht-stören-Modus: Zurückeroberung der von Benachrichtigungen geraubten Aufmerksamkeit}

\begin{itemize}
  \item \textbf{Kerngedanke}: Nicht „weniger aufs Smartphone schauen", sondern „fokussieren, wenn es Zeit zu fokussieren ist; Kontakt aufnehmen, wenn es Zeit zu kontaktieren ist".
  \item \textbf{Verwendungsempfehlungen}:\newline · Arbeitsmodus: Nur „markierte" Kontakte und notwendige Apps durchlassen;\newline · Schlafmodus: Automatisch aktivieren, nur Wecker erlauben;\newline · Sport / Lesen / Meeting: Nach aktueller Aktivität anpassen.
  \item \textbf{Fortgeschrittene Nutzung}: Mit Automatisierung + Standort-/Zeit-Trigger den Modus automatisch wechseln. Beispiel: Mit iOS-Kurzbefehlen lässt sich einstellen „bei Ankunft im Büro automatisch in den Arbeitsmodus wechseln".
\end{itemize}

\booksubsection{Benachrichtigungsverwaltung: Schluss mit dem endlosen „rote-Punkte-Ansturm“}

\begin{itemize}
  \item \textbf{Marketing-Benachrichtigungen deaktivieren}: 90 \% der Push-Nachrichten brauchst du eigentlich nicht.
  \item \textbf{Nach Wichtigkeit sortieren}: Persönliche Kommunikation (Anrufe, SMS, Instant Messaging) oben platzieren; Nachrichten, Shopping, Social Media können auf „stille Benachrichtigung" gestellt werden.
  \item \textbf{Benachrichtigungen gruppieren}: Gleichartige Benachrichtigungen zusammenfassen, um zu vermeiden, dass beim einmaligen Entsperren Dutzende Nachrichten überfluten.
\end{itemize}

\booksubsection{Schnelle Gesten: Einen Schritt schneller sein}

\begin{itemize}
  \item \textbf{Screenshot}:
  \item \textbf{Langes Screenshot / Scroll-Screenshot}: Nach einem normalen Screenshot auf die Option „Langer Screenshot" tippen.
  \item \textbf{Ein-Tasten-Scan}: Im Kontrollzentrum / in den Schnelleinstellungen ein „Scannen"-Symbol hinzufügen.
  \item \textbf{Auf Rückseite tippen (iPhone)}: Unter [Einstellungen] > [Bedienungshilfen] > [Berührung] > [Auf Rückseite tippen] Aktionen wie „Screenshot" oder „Sperren" zuweisen.
  \item \textbf{Sprachaktivierung}: Siri / Xiao Ai / Xiao Yi / Xiao Bu – ein Wort genügt, um Wetter abzufragen oder einen Wecker zu stellen.
\end{itemize}

\booksubsection{Eingabemethode: Das unterschätzte Effizienzwunder}

\begin{itemize}
  \item \textbf{Die richtige Eingabemethode wählen}:\newline · Sogou / Baidu: Große Wortdatenbank, schnelle Aktualisierung von Slang;\newline · iFlytek: Sehr hohe Spracherkennungsrate, hervorragend für Notizen;\newline · Gboard (Google): Aufgeräumt, starke plattformübergreifende Synchronisation, unterstützt Wisching-Eingabe;\newline · iOS-eigene: Bester Datenschutz, tief in das System integriert.
  \item \textbf{Must-have-Funktionen}: Intelligente Korrektur, Wortvorschläge, Cloud-Wortlisten-Synchronisation, Hochwischen für Zahlen/Symbole.
  \item \textbf{Datenschutz-Hinweis}: Vorsicht bei „Ganzer-Satz-Vorschlag" und „personalisierte Empfehlungen" – könnten Eingabeinhalte hochladen.
\end{itemize}

\booksubsection{Startbildschirm und Multitasking: Den Workflow glattziehen}

\begin{itemize}
  \item \textbf{App-Split-Screen / Bild-in-Bild}: Beim Videoschauen Nachrichten beantworten, beim Surfen Notizen machen.
  \item \textbf{Kleines Fenster / Schwebendes Fenster}: WeChat / TikTok unterstützen den „kleinen-Fenster-Modus" – nach dem Antworten sofort einklappen.
  \item \textbf{Häufige Apps anpinnen}:
\end{itemize}

\begin{tcbinfo}{„Erster Arbeitstag"-Checkliste}
Bei einem neuen Smartphone oder nach einem großen Versions-Update empfiehlt es sich, 10 Minuten in diese Einstellungen zu investieren:\newline · Sperrbildschirm-Passwort + biometrische Entsperrung · Benachrichtigungsgruppierung + Marketing-Push deaktivieren · Schriftgröße / Aktualisierungsrate / Dunkelmodus\newline · Standard-Eingabemethode / Browser / E-Mail · Notfallkontakte / medizinische ID · Backup-Einstellungen

\end{tcbinfo}

\booksection{Fünf. Bedienungshilfen und Barrierefreiheit: Die Kraft der Technik für das Gute}

Lass dich vom Namen „Bedienungshilfen“ nicht abschrecken – diese Funktionen kann jeder nutzen. Sie wurden ursprünglich für Menschen mit Behinderungen entwickelt, sind aber auch für alle anderen nützlich – etwa bei temporär verletzter Hand, beim Tragen eines Kindes, wenn im U-Bahn-Waggon ein Kopfhörer ausfällt oder wenn man mit Lesebrille kleine Schrift liest … Sie können immer helfen. Letztlich spiegelt Barrierefreiheits-Design die Idee der\textbf{Technik für das Gute}wider: Jeder soll digital gleichberechtigt teilhaben können.

\booksubsection{Die „Augen“ für sehbehinderte Nutzer – Screenreader-Komplettleitfaden}

Ein Screenreader ist das wichtigste Werkzeug für sehbehinderte Nutzer. Er liest\textbf{alle sichtbaren Bildschirminhalte vor}; der Nutzer bedient das Smartphone über Gesten und Sprache. Im Folgenden stellen wir die Verwendung von iOS „VoiceOver“ und Android „TalkBack” vor.

\begin{tcbinfo}{iOS VoiceOver}
\textbf{Aktivierungspfad}: [Einstellungen] > [Bedienungshilfen] > [VoiceOver], Schalter umlegen. Alternativ dreifaches Drücken der Seitentaste (oder Home-Taste) zum schnellen Ein-/Ausschalten.

\textbf{Grundlegende Gesten}:

\begin{enumerate}
  \item Ein-Finger-Wischen: Element für Element über den Bildschirm navigieren;
  \item Zwei-Finger-Wischen: Seite nach oben/unten/links/rechts scrollen;
  \item Ein-Finger-Doppeltipp: Aktuelles Fokuselement aktivieren (entspricht „Tippen");
  \item Drei-Finger-Wischen: Seite wechseln;
  \item Ein-Finger-Dreifach-Tipp: Vergrößern oder Auswählen (anders als normales Dreifach-Tippen).
\end{enumerate}

\textbf{Personalisierung}: [VoiceOver] > [Sprache], hier lassen sich\textbf{Sprechgeschwindigkeit}(Standard ca. 50 \%, auf 70–80 \% für flüssigeres Höhen einstellen),\textbf{Tonhöhe}einstellen und verschiedene Sprachpakete wechseln. iOS bringt „Verstärkte Sprache“ und „Sprachregler“ mit – Ersteres entfernt Atemgeräusche und Pausen, so klingt es natürlicher.

\end{tcbinfo}

\begin{tcbinfo}{: [TalkBack] > [Text-in-Sprache-Einstellungen], hier lassen sich}
\textbf{Aktivierungsweg}: [Einstellungen] > [Bedienungshilfen] > [TalkBack], Schalter aktivieren. Oder beide Lautstärketasten 3 Sekunden gedrückt halten (je nach Marke).

\textbf{Grundgesten}:

\begin{enumerate}
  \item Einfingerig links/rechts wischen: Fokus wechseln;
  \item einfingerig doppelt tippen: Element aktivieren;
  \item zweifingerig hoch/runter oder links/rechts wischen: scrollen;
  \item dreifingerig einmal tippen: in den Lesemodus wechseln (am Stück vorlesen);
  \item vierfingerig einmal tippen: Vorlesen pausieren/fortsetzen.
\end{enumerate}

\textbf{Personalisierung}: Unter [TalkBack] > [Text-zu-Sprache-Einstellungen] lassen sich\textbf{Sprechgeschwindigkeit},\textbf{Tonhöhe}und der\textbf{Sprach-Engine}wechseln (Empfehlung: Google Text-to-Speech-Engine für natürlichere Stimme).

\end{tcbinfo}

\booksubsection{Kleine Schrift groß, Unscharf scharf – Seh-Anpassungen}

\booksubsection{:}

\begin{itemize}
  \item \textbf{iOS}: [Einstellungen] > [Bedienungshilfen] > [Zoom], aktivieren, dann\textbf{drei-Finger-Tippen}auf den Bildschirm zum Vergrößern. Standardmäßig Vollbild-Zoom, umschaltbar auf Fenster-Zoom (schwebendes Vergrößerungsfenster).\textbf{Mit zwei Fingern zusammen/auseinander Zoomfaktor einstellen (1,5-fach bis 15-fach)}, drei-Finger-Doppeltipp fixiert die aktuelle Vergrößerung.
  \item \textbf{Android}: [Einstellungen] > [Bedienungshilfen] > [Vergrößerung], drei-Finger-Tippen zum Vergrößern, nochmaliges drei-Finger-Tippen zum Beenden; oder direkt\textbf{,}für temporäre Vergrößerung (maximal 8-fach).
\end{itemize}

\booksubsection{Text mit hohem Kontrast / Fetter Text}

\begin{itemize}
  \item \textbf{iOS}: [Einstellungen] > [Bedienungshilfen] > [Anzeige \& Textgröße], Schalter für „Fetter Text", „Kontrast erhöhen", „Transparenz reduzieren" usw. aktivieren. „Smart Invert" lässt Text auf hellem Hintergrund schwarz, um weißen Text auf weißem Grund zu vermeiden.
  \item \textbf{Android}: [Einstellungen] > [Anzeige] > [Schriftgröße], vier Stufen einstellbar: „Klein", „Normal", „Groß", „Sehr groß". Huawei und Xiaomi bieten zusätzlich „Fetter Text"- und „Text mit hohem Kontrast"-Schalter.
\end{itemize}

\booksubsection{iOS}

\begin{itemize}
  \item \textbf{iOS}: [Einstellungen] > [Bedienungshilfen] > [Anzeige \& Textgröße] > [Farbfilter], bietet\textbf{vier Farbblindheitsmodi (Rot/Grün/Blau)}sowie „Graustufen" und „Farben invertieren". Vorab einen Sehtest (im System integriert) durchführen und dann den passenden Filter wählen.
  \item \textbf{Android}\textbf{drei Farbblindheitsmodi}+ „Farben invertieren"-Schalter.
\end{itemize}

\booksubsection{KI-Bildbeschreibung: Die KI als deine „Augen“}

Wenn ein Screenreader auf ein Bild stößt, wird es traditionell übersprungen oder nur das Wort „Bild“ vorgelesen. Doch heute kann eine KI\textbf{iOS}– extrem hilfreich beim Erkennen von Captchas, beim Vorlesen von Diagrammen oder beim Identifizieren von Menüs.

\begin{tcbinfo}{Android}
\begin{itemize}
  \item \textbf{iOS 15+}: Unter [Einstellungen] > [Bedienungshilfen] > [VoiceOver] > [Sprachbeschreibung] „Bildbeschreibung" aktivieren. Das System nutzt ein\textbf{iOS}zur automatischen Erstellung von Bildbeschreibungen (ab iPhone 12).
  \item \textbf{: Google TalkBack integriert in neueren Versionen eine}: Google TalkBack ist in neueren Versionen integriert\textbf{Bildfrage-Feature}, mit dem du fragen kannst „Was ist auf diesem Bild?“ – die KI antwortet in Echtzeit (nicht strikt an die Android-Version gebunden; die TalkBack-App aktualisieren genügt).
  \item \textbf{Drittanbieter-Apps}: „Seeing AI" (iOS), „Lookout" (Android) und „Xiao Yi Bang Bang" (Huawei) sind hervorragende KI-Sehassistenten, die Objekte erkennen, Dokumente vorlesen und Gesichtsausdrücke identifizieren können.
\end{itemize}

\end{tcbinfo}

\begin{tcbpractice}{Anwendungsszenario}
Versuche, bei einem Straßenfoto zu fragen „Wie viele Personen sind auf diesem Bild? Was machen sie?“ oder öffne eine Lieferanten-App und lass die KI die Gerichtnamen im Menü vorlesen –\textbf{das ist wirklich „Technik, die das Leben gerechter macht”}.

\end{tcbpractice}

\booksubsection{Die „Ohren“ für hörgeschädigte Nutzer – Live-Captions-Komplettleitfaden}

Live-Captions (Echtzeituntertitel) sind eine\textbf{systemweite Sprach-zu-Text-Funktion}, die für hörgeschädigte Nutzer in lauter Umgebung, in der Bibliothek oder bei fremdsprachigen Inhalten besonders nützlich ist.

\begin{tcbinfo}{Live-Captions aktivieren und optimieren}
\textbf{iOS}: [Einstellungen] > [Bedienungshilfen] > [Live-Captions], Schalter umlegen. Das System erkennt alle Audioquellen des Geräts (Anrufe, Videos, Podcasts, FaceTime, Sprachmemos) und zeigt automatisch Untertitel oben auf dem Bildschirm an.

\textbf{Android 12+}:\textbf{Lautstärketaste}drücken und in der Popup-Anzeige auf das „Live-Captions“-Symbol tippen; oder unter [Einstellungen] > [Bedienungshilfen] > [Live-Captions] aktivieren.

\textbf{Optimierung der Genauigkeit}:

\begin{enumerate}
  \item Nach Möglichkeit in einer\textbf{.}verwenden, fern von Klimaanlagen und Maschinenlärm;
  \item Bei\textbf{deutlich artikulierter, mittlerer Sprechgeschwindigkeit}ist die Erkennungsrate am höchsten;
  \item Mehrpersonengespräche verringern die Genauigkeit – der Modus „Nur eine Person anzeigen" kann helfen;
  \item iOS\textbf{kleinere Sprachen werden schlechter erkannt};
  \item Hintergrundmusik oder Stimmengewirr können die Erkennung unterbrechen.
\end{enumerate}

\end{tcbinfo}

\booksubsection{Anruf-Szenario mit zwei Kanälen: Sprach-zu-Text + Text-zu-Sprache}

Wenn ein Nutzer\textbf{weder hören noch sprechen kann}:

\begin{tcbinfo}{Echtzeit-Anruf-Übersetzung (TTY / RTT)}
\textbf{iOS}: [Einstellungen] > [Bedienungshilfen] > [RTT/TTY], nach der Aktivierung wird beim Telefonieren eine Tastatur eingeblendet. Die Sprache des Gesprächspartners wird in Echtzeit in Text umgewandelt,\textbf{und dein getippter Text wird dem Gesprächspartner vorgelesen}.

\textbf{Android}: [Einstellungen] > [Bedienungshilfen] > [Live-Transkription] + [Ton-Benachrichtigungen]. „Live-Transkription“ wandelt die Sprache des Gesprächspartners in Echtzeit in Text um; „Ton-Benachrichtigungen” wandelt deine Eingabe in Sprache um, wenn Sprechen unpraktisch ist.

\textbf{unterstützt gängige Sprachen wie Hochchinesisch, Kantonesisch und Englisch,}:

\begin{itemize}
  \item \textbf{RTT (Real-Time Text)}: Beide Seiten müssen es unterstützen; der Text erscheint bereits beim Tippen ohne Verzögerung – die großen Netzbetreiber unterstützen es bereits.
  \item \textbf{TTY (Teletypewriter)}: Beide Seiten müssen RTT unterstützen; Text wird während des Tippens ohne Verzögerung angezeigt. Von den meisten Netzbetreibern bereits unterstützt.
  \item Die\textbf{Sprache-zu-Text-Funktion}von WeChat/QQ ist eine Alternative: Sprachnachricht lange drücken „In Text umwandeln“.
\end{itemize}

\end{tcbinfo}

\booksubsection{„Klein, aber fein“-Funktionen für alle}

\begin{itemize}
  \item \textbf{Rückseite antippen}(iOS 14+): [Einstellungen] > [Bedienungshilfen] > [Berührung] > [Auf Rückseite tippen],\textbf{Doppel- oder Dreifach-Tipp}löst Aktionen aus: Screenshot, Display sperren, Kontrollzentrum und mehr. Praktisch, wenn beide Hände ein Paket halten.
  \item \textbf{AssistiveTouch / schwebende Kugel}: iOS [Bedienungshilfen] > [Touch] > [AssistiveTouch]; Android [Bedienungshilfen] > [Schwebende Navigation]. Virtuelle Tasten ersetzen Gesten oder physische Tasten – Rettung bei defekten Tasten oder Einhandbedienung.
  \item \textbf{Schaltersteuerung}: Das Handy per externem Schalter, Einzeltipp oder Kopfbewegung steuern,\textbf{entwickelt für Nutzer mit körperlichen Einschränkungen}. iOS [Bedienungshilfen] > [Schaltersteuerung]; Android [Bedienungshilfen] > [Schalterzugriff].
  \item \textbf{Sounderkennung}(iOS): [Bedienungshilfen] > [Sounderkennung] – aktiviert erkennt das Handy Klingel, Babyweinen, Feueralarm, Sirenen und Co.,\textbf{erkennt es Vibration plus Bildschirmbenachrichtigung}. Sehr wertvoll für hörgeschädigte Nutzer und während des Schlafs.
  \item \textbf{LED-Blinkbenachrichtigung}: iOS [Bedienungshilfen] > [Audio/Visuell] > [LED-Blinken bei Hinweisen]; Android [Einstellungen] > [Benachrichtigungen] > [Blitzlicht-Hinweis]. Bei Anrufen und Benachrichtigungen blitzt die LED,\textbf{auch ohne Ton in lauter Umgebung wahrnehmbar}.
  \item \textbf{Monophonie (Ein-Kanal-Audio)}: [Einstellungen] > [Bedienungshilfen] > [Audio] > [Mono-Audio] vereint linken und rechten Kanal. Wer auf einem Ohr schwer hört, verliert so keine Hälfte der Information.
\end{itemize}

\begin{tcbinfo}{Apples Barrierefreiheits-Philosophie}
Apple investiert viel in Barrierefreiheit und erweitert die Funktionen jedes Jahr. Auf der Apple-Website gibt es eine eigene Barrierefreiheitsseite mit detaillierten Anleitungen. iOS ist in diesem Bereich meist ausgereifter und flüssiger als Android – ein wichtiger Grund, warum viele Nutzer mit Behinderung Apple wählen.

\end{tcbinfo}

\begin{tcbpractice}{30 Sekunden ausprobieren}
Nimm dir etwas freie Zeit für folgende Erfahrungen:

\begin{enumerate}
  \item Aktiviere\textbf{VoiceOver / TalkBack}, versuche mit drei Fingern durch WeChat zu blättern – ganz ohne Augen;
  \item aktiviere die\textbf{Live-Untertitel}, sprich aufs Handy und prüfe die Erkennungsquote;
  \item nutze den\textbf{dreifingerigen Doppeltipp}einen Text auf 200 \%, um den Alltag sehbehinderter Nutzer zu erleben;
  \item Deaktiviere alle zuvor aktivierten Funktionen und überlege: „Wie würde ich sie gestalten wollen, wenn ich sie dauerhaft bräuchte?"
\end{enumerate}

\end{tcbpractice}
